\section{Strengths and Limitations}

The main strength of this route is that it does not confuse evaluation with decision. Many weak drafts stop at a comprehensive ranking table and implicitly treat the top-ranked supplier as the final answer. The present route avoids that shortcut by giving the evaluation model a narrower but more defensible role: it reduces the candidate space and explains why certain suppliers deserve to enter the planning stage. The actual recommendation is then produced by a separate decision layer with explicit feasibility constraints.

The second strength is that the evidence chain is visible. The paper can point to normalized indicators in the data section, score and retained-candidate artifacts in the result section, an executable allocation table in the planning layer, and scenario or feasibility checks in the validation section. This makes the paper read more like a completed team draft than a prestige-method outline, because each named method closes into a concrete artifact and an interpretable conclusion.

The limitations are equally important to state. First, the evaluation layer depends on the quality of the indicator set and weighting logic; if the indicators omit important operational burden, then the retained candidate set may already be biased before the planning model begins. Second, the planning model usually simplifies procurement reality by compressing lead time, batch rules, dynamic supplier interactions, or multi-material coupling into a smaller feasibility system. Third, any scenario analysis remains conditional on the chosen disturbance range; a stable recommendation under mild perturbation does not automatically imply resilience under extreme disruptions.

For later iterations, two extensions are especially valuable. One is to strengthen the comparison layer by adding baseline methods or alternative candidate-screening strategies, so that the retained supplier set is not justified by only one score mechanism. The other is to deepen the executable planning layer through richer constraints, multi-period inventory linkage, or more realistic uncertainty propagation. These extensions preserve the same paper route while making the final recommendation more competitive.
